\section*{元数学}
\phantomsection
\addcontentsline{toc}{section}{元数学}

无矛盾性在任何时代都被视为全部数学的一个必要条件，而且从亚里士多德的时代起，逻辑就已充分发展，人们完全清楚地知道：从一个矛盾的理论可以推出任何东西。自古代以来就被视为不可或缺的那些“存在性”证明，显然别无其他目的，只是要保证新概念的引入不至于带来矛盾，尤其是当这个概念过于复杂、无法直接纳入“直观”之下时。我们已经看到，随着 19 世纪公理学观点的到来，这一要求变得如何迫切，而算术“模型”的构造又是如何回应它的。但是，算术本身难道不会是矛盾的吗？在 19 世纪末之前，人们无疑不会想到这样的问题，因为整数显得属于我们直观中最可靠的那一部分；但在“悖论”之后，一切都似乎成了问题，于是可以理解，它们所造成的不安全感促使数学家们从 1900 年前后起更加慎重地对待算术的无矛盾性问题，以便至少使经典数学免遭灭顶之灾。这个问题因此成为希尔伯特在其著名的 1900 年国际数学家大会演讲中所列举的第二个问题（{[}163 a{]}，第 III 卷，第~229-301~页）。在这样做时，他提出了一条将产生巨大反响的新原则：在传统逻辑中，一个概念的无矛盾性只是使它成为“可能的”，而对希尔伯特来说（至少对公理地定义的数学概念而言），无矛盾性与概念的存在是一回事。这显然意味着，必须先证明一门数学理论的无矛盾性，然后才能合法地发展它；H. 庞加莱正是这样理解的，当他在形式主义身上打开缺口时，他把希尔伯特的这一想法接了过来，并怀着狡黠的快意强调，形式主义者在当时离能够做到这一点还相去甚远（{[}251 c{]}，第~163~页）。我们将在后文看到希尔伯特如何应战；但在那之前，这里必须指出，在这两位的权威的影响下，庞加莱所提出的这些要求在很长时期内被形式主义者及其论敌双方毫无保留地接受了。其后果之一是形式主义者中流传甚广的一个信念：希尔伯特的证明论是数学的一个组成部分，而且是其不可或缺的绪论。这一教条在我们看来并无根据，\(^{67}\) 我们认为，元数学对逻辑和数学的阐述的介入，可以而且应当缩减到涉及缩写记号的运用和演绎判据的那一点非常初等的部分。因此，与庞加莱的想法相反，问题不在于“要求免于矛盾”，而在于同阿达马一起认为：无矛盾性即使未被证明，也是被注意到的（{[}12{]}，第~270~页）。

剩下要做的，是对希尔伯特及其学派的努力作一简短的历史概述，并且概略地追述两条线索：一是最终导出哥德尔否定性结果、并使阿达马的怀疑态度得到事后印证的那段演变，二是此后围绕数学推理机制的认识所取得的一切进步——正是这些进步使现代元数学成为一门自有其不容置疑的意义的独立科学。

早在 1904 年，希尔伯特就在国际数学家大会的一次演讲（{[}163 c{]}，第~247-261~页）中着手处理算术的一致性问题。他首先指出，利用模型来证明它是不可能的，\(^{68}\) 并概述了另一种方法的原理：他提议把形式化算术的真命题看作没有意义的符号汇集，并证明：在使用支配这些汇集的构成与联结的规则时，不可能得到一个既是真命题、其否定又是真命题的汇集。他甚至对一个比算术更窄的形式勾勒了这种性质的证明；但是，正如 H. 庞加莱不久之后所指出的（{[}251 c{]}，第~185~页），这个证明本质上使用了递归原理，因而看来依赖于一个恶性循环。希尔伯特没有立即回应这一批评，此后大约十五年无人尝试发展他的思想；直到 1917 年（出于回应直觉主义者攻击的愿望），他才重新套上数学基础问题这驾马车，而且从此直到科学生涯结束不曾停止考虑这个问题。在他大约从 1920 年延续到 1930 年的关于这一主题的工作中——一整派年轻数学家（阿克曼、贝尔奈斯、埃尔布朗、冯·诺伊曼）参与了这些工作——希尔伯特一点一点地更精确地提出了他的“证明论”的诸原则：他默认了庞加莱批评的正确性，承认在元数学中，所使用的算术推理只能建立在我们对整数的直观（而不是形式化算术）之上；为此，看来必须把这种推理限制在直觉主义者所容许的那类“有限程序”（“finite Prozesse”）之内：例如，一个反证法证明不能证明一个汇集或汇集序列在元数学上的存在，必须有一个明确的构造规律。\(^{69}\) 另一方面，希尔伯特在两个方向上扩大了他最初的纲领：他不仅处理算术的一致性，还力图证明实数论甚至集合论的一致性；\(^{70}\) 此外，在一致性问题之外又添上了公理的独立性、完备性和判定问题。我们将逐一简述这些不同的问题，并指出它们所引出的主要研究。

命题组 \(A_{1}, A_{2}, \ldots, A_{n}\) 的独立性证明在于表明：对每个指标 i，\(A_{i}\) 在理论 \(T_{i}\) 中不是定理，该理论是以 \(A_{j}\)（其指标 \(j \neq i\)）为公理而得到的。如果知道一个一致的理论 \(T_{i}^{\prime}\)，在其中 \(A_{j}(j \neq i)\) 都是定理，而且“非 \(A_{i}\)”也是定理，这一点便得到了确立；于是，根据是否承认某些理论（如算术或集合论）是一致的，可以以两种方式考虑这个问题。在第二种情形，所面对的是一个“绝对”一致性问题。相反，第一类问题通过构造适当的“模型”化归为“相对”一致性问题，而且早在数学取得完全形式化的面貌之前，人们就已想出了许多这种类型的证明：只需回想非欧几何的模型（见第~133~页）、希尔伯特在《几何基础》（Grundlagen der Geometrie）中处理的初等几何公理的独立性问题 {[}163 c{]}，以及施泰尼茨关于代数公理化的工作和豪斯多夫及其后继者关于拓扑学公理化的工作，即足以说明。

称一个理论 T 为完备的，如果对 T 的每个命题 A（不含 T 的常元以外的任何字母），命题 A 与（非 A）之一是 \(T.^{71}\) 的定理。若把某些非常初等的形式体系撇在一旁——它们的完备性是容易证明的（{[}164{]}，第~35~页）——这个领域所得的结果本质上是否定性的；其中最早的一个属于 K. 哥德尔 {[}130 a{]}，他表明：如果 T 是一致的，且形式化算术的公理都是 T 的定理，那么 T 不是完备的。他那个巧妙方法的根本想法，在于（当然借助于“有限程序”）在元数学陈述与形式化算术中的某些命题之间建立一个双射对应；我们只限于勾画一个粗略的轮廓。\(^{72}\) 对 T 的每个作为项或关系的汇集 A，先以一个可以近乎机械地应用的（属于明确构造的）程序，以双射的方式关联一个整数 \(g(A)\)。同样，对 T 的每个证明 D（被看作汇集的一个相继序列），可以以双射的方式关联一个整数 \(h(D)\)。最后，可以给出一个明确地构造 T 的关系 \(P(x,y,z)\) 的程序，\(^{73}\) 使得在 T 中，\(P(x,y,z)\) 要求 x,y,z 是整数，并且满足以下两个条件：

\(1^{\circ}\) 若 \(D\) 是 \(A(\lambda)\) 的一个证明，其中 \(A(x)\) 是 \(\mathcal{T}\) 的一个关系，而 \(\lambda\) 是一个明确的整数（也就是说，是 \(\mathcal{T}\) 的一个是整数的项），则 \(P(\lambda, g(A(x)), h(D))\) 是 \(\mathcal{T}\) 的定理；

\(2^{\circ}\) 若明确的整数 \(\mu\) 不具有 \(h(D)\) 的形式，或者 \(\mu = h(D)\) 而 \(D\) 不是 \(A(\lambda)\) 的证明，则（非 \(P(\lambda, g(A(x)), \mu)\)）是 \(\mathcal{T}\) 的定理。

现在设 \(S(x)\) 是关系（非 \((\exists z)P(x,y,z)\)），并设 \(\gamma = g(S(x))\)，它是 \(\mathcal{T}\) 的一个项。如果 \(\mathcal{T}\) 是一致的，那么命题 \(S(\gamma)\) 在 \(\mathcal{T}\) 中就没有证明。事实上，若 \(D\) 是这样一个证明，则 \(P(\gamma,g(S(x)),h(D))\) 将是 \(\mathcal{T}\) 的定理；但这个关系不是别的，正是 \(P(\gamma,\gamma,h(D))\)，由此推知 \((\exists z)(P(\gamma,\gamma,z)\) 也将是 \(\mathcal{T}\) 的定理；由于这最后一个关系等价于（非 \(S(\gamma)\)），\(\mathcal{T}\) 就会是不一致的。另一方面，刚才所说的表明，对每个明确的整数 \(\mu\)，（非 \(P(\gamma,\gamma,\mu)\)）都是 \(\mathcal{T}\) 的定理。其结果是，\(\mathcal{T}\) 中不存在（非 \(S(\gamma)\)）的证明，因为这最后一个关系等价于 \((\exists z)P(\gamma,\gamma,z)\)，而一旦存在整数 \(\mu\) 使 \(P(\gamma,\gamma,\mu)\) 成立，依据前文所述就将要求 \(\mathcal{T}\) 是不一致的。\(^{74}\) 哥德尔的这条元数学定理随后在若干方向上得到了推广（{[}181{]}，第 XI 章）。\(^{75}\)

T 的关系 \(S(\gamma)\)——我们刚才表明，无论这个关系还是它的否定，在 T 中都没有证明——显然是专为达成目的而拼造出来的，它同任何数学问题都没有自然的联系。有意思得多的是如下事实：若以 \(\mathcal{T}\) 表示集合论（采用冯·诺伊曼-贝尔奈斯公理系统），则连续统假设及其否定在 \(\mathcal{T}\) 中都不可证明。这一著名结果是分两步建立的：1940 年，哥德尔证明了通过向 \(\mathcal{T}\) 添加连续统假设 \(2^{\aleph_0} = \aleph_1\) 而得到的理论并非不一致 {[}130 b{]}；随后在 1963 年，P. 科恩证明了：当向 \(\mathcal{T}\) 添加关系 \(2^{\aleph_0} = \aleph_2\)（或者，对任意整数取 \(2^{\aleph_0} = \aleph_n\)，其中 \(n > 1\)）时，结论同样成立 {[}67{]}。

判定问题（“Entscheidungsproblem”）无疑是元数学中所陈述的一切问题里雄心最大的一个：它在于弄清，对一个给定的形式语言，能否设想一种近乎机械的“普适程序”，当把它应用于所论形式体系的任何关系时，能在有限次运算之内指出该关系是否为真。\(^{76}\) 这一问题的解决本来就属于莱布尼茨的宏伟设想的一部分，而且希尔伯特学派有一度似乎认为其实现已近在眼前。事实上，对只含少数原始符号和公理的形式体系，人们能够描述出这样的程序（{[}181{]}，第~136-141~页）。但是，为使判定问题精确化、即确切勾勒“普适程序”意指什么而作的种种努力，迄今为止只得到了否定的结果（{[}181{]}，第~432-439~页）。此外，一个理论 T 的判定问题的解决，立即使人得知 T 是否一致，因为只需把“普适程序”应用于 T 的一个关系及其否定即可；而我们将看到，对通常的数学理论，不可能以这种方式解决这一问题。\(^{77}\)

的确，正是在数学理论的一致性问题——元数学的源头和心脏——上，结果显得最令人失望。在 1920-1930 年间，希尔伯特及其学派为攻克这些问题发展了新的方法；在证明了覆盖算术之一部分的部分形式体系的一致性之后（参见 {[}158{]}、{[}324 d{]}），他们以为自己已经达到目标，证明了不仅是算术的一致性、而且还有集合论的一致性；这时哥德尔依据算术的不完备性推出：用希尔伯特的“有限程序”不可能证明任何包含算术的理论 T 的一致性。\(^{78}\)

然而，哥德尔定理并没有把证明一致性的尝试之门完全关死，只要（至少部分地）放弃希尔伯特关于“有限程序”的限制即可。于是根岑在 1936 年 {[}127{]} 成功地证明了形式化算术的一致性，他所使用的是“直观地”超限归纳到可数序数 \(\varepsilon_{0}\) 的做法。⁷⁹ 可以加于这种推理之上的“确定性”的价值，无疑不如满足希尔伯特最初要求的推理那样令人信服，而且它本质上是每个数学家的个人心理问题；但同样真实的是：这类“直观地”超限归纳到某一给定序数的类似“证明”，若能应用于例如实数理论或集合论的一个实质部分，将会被视为重要的进步。

